% Generated draft: requires mathematical and source review.
\documentclass[11pt]{article}
\usepackage[margin=25mm]{geometry}
\usepackage{fontspec}
\setmainfont{Latin Modern Roman}
\usepackage{amsmath,amssymb,mathtools,unicode-math}
\setmathfont{Latin Modern Math}
\usepackage{enumitem,needspace,xurl,hyperref}
\setlength{\parindent}{0pt}
\setlength{\parskip}{5pt}
\setlength{\emergencystretch}{3em}
\newcommand{\N}{\mathbb{N}}
\newcommand{\Z}{\mathbb{Z}}
\newcommand{\Q}{\mathbb{Q}}
\newcommand{\R}{\mathbb{R}}
\newcommand{\C}{\mathbb{C}}
\newcommand{\F}{\mathbb{F}}
\newcommand{\A}{\mathbb{A}}
\newcommand{\PP}{\mathbb{P}}
\newcommand{\E}{\mathbb{E}}
\newcommand{\eps}{\varepsilon}
\newcommand{\dd}{\,\mathrm{d}}
\begin{document}
\section*{UnsolvedMath ID: 30000338}
\textbf{Source problem code:} catalogue entry
\subsection{Definitions and mathematical statement}
The source asks whether model selection means choosing a statistical model for the data, choosing a prediction function, or averaging either object. This is an objective-setting question, not a supplied assertion that one interpretation is universally correct or that the alternatives are mutually exclusive.

To distinguish the targets, let $D$ denote observed data and $\{\mathcal P_m:m\in\mathcal I\}$ a collection of statistical models, each a family of probability laws. Model selection returns an index $\widehat m(D)$, potentially followed by a fitted law in $\mathcal P_{\widehat m}$. Prediction selection instead returns a rule $\widehat f_D$ judged by a chosen future-observation risk, for example
\[
 R(f)=\mathbb E[L(Y,f(X))],
\]
where the joint law of $(X,Y)$ and the loss $L$ must be specified. Choosing the data-generating model and minimizing prediction risk are different criteria.

For a finite collection of fitted laws $P_{m,D}$ and predictions $f_{m,D}$, data-dependent weights $w_m(D)\geq0$, $\sum_mw_m(D)=1$, yield respectively
\[
 P_D^{\mathrm{av}}=\sum_m w_m(D)P_{m,D},\qquad
 f_D^{\mathrm{av}}=\sum_m w_m(D)f_{m,D},
\]
where the second expression requires an output space in which averaging is defined. Mixture distributions and averaged point predictions are different objects; the latter need not coincide with the prediction induced from the former for a nonlinear loss or decision rule. The source does not choose Bayesian versus frequentist weights, a loss, a model class, a true-model assumption or an asymptotic criterion.

The research task is to clarify which target and assessment criterion fit the intended inferential problem, and how their objectives differ. The literature triage checked 21 August 2026 explicitly describes this as a methodological question rather than a theorem with a settled universal answer. The registry's open label is retained in that qualified sense.
\subsection{Short English statement}
The entry asks which object model selection should choose: a model for the data, a prediction rule, or an average. These correspond to different objectives and risks. Without a specified inferential task and loss, the source does not determine one universally preferred target.
\subsection{Sources}
\begin{itemize}
\item Local UnsolvedMath registry, permanent ID 30000338, exact source code catalogue entry; \emph{Statistical Model Selection and Prediction Targets}. Proposed year 2005, status open; exact methodological wording retained.
\item Supplied report citation: \emph{Statistische und Probabilistische Methoden der Modellwahl (2006)}, Oberwolfach Reports, DOI 10.4171/owr/2005/47, pp. 2611--2704, \url{https://doi.org/10.4171/owr/2005/47}. Title year differs from DOI/proposed year.
\item Supplied literature triage checked 2026-08-21, citing only ``original Oberwolfach report'', \url{https://doi.org/10.4171/owr/2005/47}. It classifies the item as an objective-setting methodological research question. The dated OWR-derived curation requests independent verification of time-sensitive status and cited claims.
\end{itemize}
\end{document}
